% ---------------------------------------------------------------------------
%  fig2.tex -- Figure 2 of the CMOS chain paper: the tensor-network method.
%
%  Replaces NewFigures/Figure2CMOSVertical7.pdf, a PowerPoint export set in
%  Cambria Math / Arial / Aptos / DejaVuSans, whose MediaBox also carries a
%  superseded draft of the whole diagram outside the CropBox.
%
%  Authored at exactly \columnwidth = 246pt, so \includegraphics applies no
%  scaling: include with [width=\columnwidth].  The original was 204 x 427pt
%  shown at 162.6pt wide; re-proportioned here by lifting the basis-function
%  inset out of the vertical stack and setting it beside the eigenvalue-
%  equation / MPS stage, which it annotates as a zoom callout.
%
%  Data: fig2_basis.dat, written by prep_fig2.py from
%  ../PostReviewFigures/Figure 2 Data/BasisVectors{X,Y}Axis.txt
%
%  revB/C: the change of basis W -> U W U^dagger is DRAWN.  Two operators now
%  carry two colours: the REDUCED MPO (physical dimension d) is teal, the FULL
%  MPO (physical dimension D) is blue.  Three sites show only the reduced
%  tensor; the rightmost site is opened up to show what that tensor is, as a
%  teal-stroked box enclosing U^dagger, the full blue MPO, and U.  The same
%  object therefore appears at two zoom levels, which is what supplies the
%  identity without writing one out.
%
%  revD: RE-LAID OUT, not rebuilt.  Same tensors, boxes, labels and curves as
%  revC, in four bands top to bottom:
%      1. the contraction W|P_i>, with the operator note and the basis inset
%         in the right-hand column;
%      2. "= mu_i" and the result MPS, shifted right to x = 110..232;
%      3. the CMOS chain, directly BENEATH the result MPS and sharing its
%         abscissae exactly, so a chain node sits under the MPS site that
%         carries it;
%      4. the two DMRG variants and their eigenpairs, at the far left, reached
%         by arrows fanning down-left out of the equation.
%
%  Why the chain moved to the bottom.  It used to run across the top, and its
%  nodes did not sit above the matching tensor sites, so the correspondence was
%  carried by repeating v_0 ... v_L on both the chain and the result MPS.  That
%  second set of labels was wrong: after U^dagger those legs run over the d=32
%  RETAINED COEFFICIENTS, not over a node voltage.  The only legs in the figure
%  that legitimately run over voltages are the two inner sandwich legs marked
%  D=65.  Setting the chain directly under the MPS lets ALIGNMENT carry the
%  correspondence -- "this tensor is that node", true in any basis -- and the
%  duplicate labels are gone.  The chain is back on its unstretched geometry
%  (units 22pt wide on a 44/34/44 site pitch), shifted +100 to sit under the
%  result MPS; the stretched full-width version revC used is not compatible
%  with alignment, which is the stronger claim.
%
%  Orientation.  \tilde{W}_{ab} = sum_{vv'} U_{av} W_{vv'} U^dagger_{v'b} with
%  a, b the d-indices; acting on a reduced-space vector b is the INPUT index
%  and a the OUTPUT.  Here the MPS sits above the operator and the free legs
%  point down, so the input leg is the upper one: U^dagger faces the MPS and U
%  sits on the lower free leg.  Reading the site DOWNWARDS therefore spells
%  U^dagger W U, while the operator identity written beside it is
%  W -> U W U^dagger.  Both are correct; they are the same statement read in
%  opposite directions.  (An earlier version of this comment said the site
%  reads U W U^dagger.  It does not; the drawing was always right.)
%
%  The zoom callout is anchored to the lower U box, not to the MPS site,
%  because the rows of U are exactly what the inset plots.  With the inset now
%  level with the sandwich the callout is a short 17pt wedge rather than the
%  long diagonal it used to travel.
%
%  NOTE on units: this picture sets x=1pt,y=1pt so bare coordinates are
%  points, but pgfplots' own at/width/height keys must still carry explicit
%  units -- unitless values there are not read in the picture's x/y vectors.
% ---------------------------------------------------------------------------
\documentclass[border=0pt]{standalone}
% ---------------------------------------------------------------------------
%  figstyle.tex -- shared style for every figure in the CMOS chain paper.
%  \input this from each figN.tex. Do not restyle locally; change it here.
%
%  Design rule: every figure is authored at EXACTLY the width it occupies on
%  the page, so \includegraphics applies no scaling and a 9pt label in one
%  figure is a 9pt label in every other. The paper's geometry, measured from
%  revtex4-2 [reprint,aps,pre]:
%      \columnwidth = 246.0pt = 3.404in   -> \figcolwidth
%      \textwidth   = 510.0pt = 7.057in   -> \figfullwidth
%  Include single-column figures with [width=\columnwidth] and full-width
%  figures (figure*) with [width=\textwidth]. Nothing else.
%
%  Verifying the width: pdfinfo reports big points (72/in), not LaTeX pt
%  (72.27/in). A correct column figure reads 245.08 bp; a correct full-width
%  figure reads 508.09 bp. Those are the numbers to check, not 246 and 510.
%
%  Fonts: Computer Modern, i.e. the LaTeX default. Load no font package.
%  Class sizes for reference: body 10, \small 9, \footnotesize 8, \scriptsize 7.
% ---------------------------------------------------------------------------
\usepackage{tikz}
\usepackage{pgfplots}
\usepackage{amsmath}
\usepackage{braket}
\pgfplotsset{compat=1.17}
\usetikzlibrary{arrows.meta,positioning,calc,backgrounds,fit,decorations.pathreplacing}
\usepgfplotslibrary{colormaps}

% ---- the two canonical widths -------------------------------------------
\newlength{\figcolwidth}   \setlength{\figcolwidth}{246.0pt}
\newlength{\figfullwidth}  \setlength{\figfullwidth}{510.0pt}

% ---- force an exact output width ----------------------------------------
%  Put \figcanvas{<width>}{<height>} as the FIRST line inside a tikzpicture.
%  It pins the bounding box, so the emitted PDF is exactly that size no
%  matter what the content does, and no rescaling happens at include time.
\newcommand{\figcanvas}[2]{\useasboundingbox (0,0) rectangle (#1,#2);}

% ---- type sizes (true points, because scale is 1) -----------------------
\newcommand{\figaxislabel}{\small}       % 9pt  -- axis labels
\newcommand{\figticklabel}{\footnotesize}% 8pt  -- tick labels
\newcommand{\figannot}{\footnotesize}    % 8pt  -- in-figure annotation
\newcommand{\figpanel}{\small\bfseries}  % 9pt  -- panel letters (a), (b)

% ---- palette -------------------------------------------------------------
\definecolor{figblueA}{HTML}{1F4E9C}  % dark blue   -- emphasis curve 1
\definecolor{figblueB}{HTML}{3C8DBC}  % mid blue    -- emphasis curve 2
\definecolor{figgreen}{HTML}{9CC7A9}  % pale green  -- de-emphasised curves
\definecolor{figgray}{HTML}{7F7F7F}   % neutral rules and guides
\definecolor{figred}{HTML}{C0392B}    % the one warning/annotation colour
\definecolor{fignodeA}{HTML}{6E8B74}  % node dot, k=0 side
\definecolor{fignodeB}{HTML}{1F4E9C}  % node dot, k=1 side
\definecolor{figshadeA}{HTML}{EAEFEA} % k=0 panel background
\definecolor{figshadeB}{HTML}{DCE8F5} % k=1 panel background
% second operator colour: the reduced MPO (physical dimension d), as against
% the full MPO (physical dimension D) drawn in figblueA/figblueB.
\definecolor{figtealA}{HTML}{1B6E6A} % reduced-MPO outline
\definecolor{figtealB}{HTML}{9FCFCB} % reduced-MPO fill

% ---- the V_dd ramp -------------------------------------------------------
%  Four samples of viridis at 0.08, 0.36, 0.64 and 0.92, the ramp the
%  matplotlib composite used, so a given V_dd is the same colour in every
%  panel that shows one.  The values are read out of the content stream of
%  PostReviewFigures/CompositeFigure3_v3.pdf, so they are matplotlib's own.
\definecolor{figvddA}{rgb}{0.283197,0.115680,0.436115}  % V_dd/V_T = 1.1
\definecolor{figvddB}{rgb}{0.179019,0.433756,0.557430}  % V_dd/V_T = 1.2
\definecolor{figvddC}{rgb}{0.170948,0.694384,0.493803}  % V_dd/V_T = 1.3
\definecolor{figvddD}{rgb}{0.793760,0.880678,0.120005}  % V_dd/V_T = 1.4

% ---- line weights --------------------------------------------------------
\pgfplotsset{
  figframe/.style={line width=0.6pt},
  figcurve/.style={line width=1.0pt},
  figemph/.style={line width=1.4pt},
}
\tikzset{
  figpanellabel/.style={font=\small\bfseries, inner sep=1pt},
  figwire/.style={line width=0.9pt},
  figlead/.style={figgray, line width=0.5pt, densely dotted},
  figarrow/.style={-{Stealth[length=4pt,width=3pt]}, line width=0.7pt},
}

% ---- the common axis style ----------------------------------------------
%  Matches the tick conventions already used by the matplotlib figures:
%  inward ticks on all four sides, minor ticks visible, no grid.
\pgfplotsset{
  figaxis/.style={
    scale only axis,
    axis line style={line width=0.6pt},
    tick align=inside,
    tickpos=both,
    minor tick num=1,
    tick label style={font=\figticklabel},
    label style={font=\figaxislabel},
    legend style={font=\figannot, draw=none, fill=none},
    every axis plot/.append style={figcurve},
  },
  figdensity/.style={
    figaxis,
    % `viridis high res' is matplotlib's own 256-entry table (the plain
    % `viridis' key is an 18-point approximation to it), so a density panel
    % ported from matplotlib keeps the colours it had.
    colormap/viridis high res,
    point meta min=0,
  },
}

\usepackage{amssymb}   % \mathbb, as used for \mathbb{W} in the body text

\begin{document}
\begin{tikzpicture}[x=1pt,y=1pt]
% Content occupies y in [113.1,377.8] -- measured with gs -sDEVICE=bbox on a
% build whose box was opened to (0,0)-(\figcolwidth,500), then padded by ~2pt
% each way.  It sits in that band rather than starting at the origin, so the
% box is pinned explicitly instead of with \figcanvas.  Re-measure and update
% these two numbers if the vertical extent of anything changes.  Result:
% 246 x 269pt, which pdfinfo reports as 245.08 x 268.00 bp.
\useasboundingbox (0,111) rectangle (\figcolwidth,380);

% ===========================================================================
%  local drawing idioms.  Geometry only: every colour, weight and type size
%  is taken from figstyle.tex, never re-declared here.
% ===========================================================================
\tikzset{
  figannot/.style={font=\figannot, inner sep=1pt},
  figemph/.style={/pgfplots/figemph},
  % MPS tensor: grey circle.
  mpsnode/.style={circle, draw=black, line width=0.9pt, fill=black!12,
                  minimum size=16pt, inner sep=0pt},
  % REDUCED MPO tensor (physical dimension d): teal rounded square.  This is
  % the object the reader meets first, at the three unexpanded sites.
  %  Drawn with the ENCLOSURE's aspect and corner radius, not as a square, so
  %  that the closed tensor and the opened one read as the same object.
  mpored/.style={rounded corners=5pt, draw=figtealA, line width=0.9pt,
                 fill=figtealB!75, minimum width=18pt, minimum height=34pt,
                 inner sep=0pt},
  % FULL MPO tensor (physical dimension D): blue rounded square.  It appears
  % only inside the expanded site, between U^dagger and U.
  mpofull/.style={rounded corners=4pt, draw=figblueA, line width=0.9pt,
                  fill=figblueB!85, minimum size=18pt, inner sep=0pt},
  % the box enclosing U^dagger / full MPO / U.  Stroked in the reduced-MPO
  % colour and filled opaquely, so it reads as the SAME object as the teal
  % squares at the other three sites -- drawn open instead of closed.
  mporedbox/.style={rounded corners=5pt, draw=figtealA, line width=0.9pt,
                    fill=figtealB!25},
  bond/.style={line width=1.4pt},              % contraction line
  % the change-of-basis isometry U: a rank-2 box, sharp-cornered, neutral grey
  % and white-filled so it reads as neither an MPS circle nor an MPO rounded
  % square.  Opaque, so it breaks the leg it sits on and leaves a visible stub
  % either side.  The fixed text height/depth keeps U and U^dagger on a common
  % baseline.
  ubox/.style={draw=figgray, line width=0.9pt, fill=white,
               minimum width=16pt, minimum height=11pt, inner sep=0pt,
               text height=5.5pt, text depth=0pt, font=\figannot},
  callout/.style={figblueB!45, line width=0.7pt},
  fatarrow/.style={-{Stealth[length=7pt,width=6pt]}, line width=2.2pt,
                   figblueA},
}

% cross-coupled inverter pair.  #1 = x of box left edge, #2 = y of the wire.
% Box 22 x 26pt; each triangle has its vertical base on one side edge and its
% apex on the opposite horizontal edge, so the pair maps onto itself under a
% 180-degree rotation about the box centre.  White fill hides the box edge
% running behind each triangle, as in the original.
\newcommand{\cmosunit}[2]{%
  \begin{scope}[shift={(#1,#2)}]
    \draw[figwire] (0,-13) rectangle (22,13);
    \draw[figwire,fill=white] (2.2,20.8) -- (2.2,5.2) -- (14.7,13) -- cycle;
    \draw[figwire,fill=white] (19.8,-20.8) -- (19.8,-5.2) -- (7.3,-13) -- cycle;
  \end{scope}}

\newcommand{\nodedot}[2]{\fill (#1,#2) circle[radius=1.9pt];}

% three-dot ellipsis centred at (#1,#2)
\newcommand{\ellip}[2]{%
  \foreach \dx in {-3.6,0,3.6}{\fill ({#1+\dx},#2) circle[radius=1.15pt];}}

% a line broken by an ellipsis: #1 style, #2 x1, #3 x2, #4 y, #5 x of the
% ellipsis.  \brokenline centres it; \brokenlineat takes it explicitly, which
% the operator row needs -- the enclosure eats into the gap from the right, so
% an ellipsis centred between the two SITES would run into its wall.
\newcommand{\brokenlineat}[5]{%
  \draw[#1] (#2,#4) -- ({#5-5.6},#4);
  \draw[#1] ({#5+5.6},#4) -- (#3,#4);
  \ellip{#5}{#4}}
\newcommand{\brokenline}[4]{%
  \pgfmathsetmacro{\bmid}{(#2+#3)/2}%
  \brokenlineat{#1}{#2}{#3}{#4}{\bmid}}

% ---- the four site abscissae of the contraction (band 1) -----------------
\def\xA{10}\def\xB{54}\def\xC{88}\def\xD{132}
\def\xcol{155}                      % left edge of the right-hand column

% ---- band 2 and 3 share ONE set of abscissae -----------------------------
%  The result MPS and the chain beneath it are the same four sites, so they
%  are drawn from the same numbers.  This is the alignment that replaced the
%  duplicated v_0 ... v_L labels; do not let the two rows drift apart.
%  The pitch is 46/34/46 rather than the tensor rows' 44/34/44: each elided
%  span has to hold a 22pt inverter unit AND an ellipsis, and at 44 the three
%  ellipsis dots ended 1.35pt from the next node dot -- the same gap they keep
%  from each other -- so the node read as a fourth omitted-unit dot.  At 46 the
%  gap is 2.85pt against a dot that is also 1.6x larger, and it separates.
%  The right-hand site is at 236: its MPS circle reaches x=244, inside 246.
\def\rA{110}\def\rB{156}\def\rC{190}\def\rD{236}

% ---- the expanded site: enclosure and the stations of its physical leg ----
%  Reading downwards: MPS (leg dimension d) -- U^dagger -- full MPO (leg
%  dimension D) -- U -- free physical leg (dimension d).  The enclosure spans
%  the three operator boxes; the MPS above it and the free leg below it stay
%  outside, because they are not part of the reduced tensor.
\def\ymps{362}\def\yUa{334}\def\ympo{310}\def\yUb{287}\def\yleg{265}
\def\encL{115}\def\encR{149}\def\encT{345}\def\encB{276}
%  The stubs of bare leg between boxes are 9-11pt, not the 5.5pt they were:
%  a dimension label set beside the leg needs a segment it can sit level with,
%  and at 8pt type that segment has to be about 9pt of clear line.
\def\xdim{128}                      % right edge of the leg-dimension labels

% ===========================================================================
%  1.  the eigenvalue problem, drawn.  The MPS (grey) contracted with the
%      reduced MPO (teal); the rightmost site opened up to show that the
%      reduced tensor is U^dagger W U.
% ===========================================================================
% The enclosure goes down FIRST, as a background: its fill is opaque, so any
% bond or leg drawn before it would be erased.  Everything the box contains --
% the virtual bond entering from the left, the physical leg running through --
% is then drawn over it, which is also how it should read: the legs pierce the
% box.  No pale highlight column here; the teal box is itself the marker for
% which site has been opened up.
\draw[mporedbox] (\encL,\encB) rectangle (\encR,\encT);
\brokenline{bond}{\xA}{\xB}{\ymps}
\draw[bond] (\xB,\ymps) -- (\xC,\ymps);
\brokenline{bond}{\xC}{\xD}{\ymps}
\brokenline{bond}{\xA}{\xB}{\ympo}
\draw[bond] (\xB,\ympo) -- (\xC,\ympo);
\brokenlineat{bond}{\xC}{\xD}{\ympo}{106}
% one continuous vertical line per site: the opaque tensors laid over it below
% carve out the MPS-MPO contraction, the U sandwich and the free physical leg
\foreach \x in {\xA,\xB,\xC,\xD}{%
  \draw[bond] (\x,\ymps) -- (\x,\yleg);
  \nodedot{\x}{\yleg}}
\foreach \x in {\xA,\xB,\xC}{%
  \node[mpsnode] at (\x,\ymps) {};
  \node[mpored]  at (\x,\ympo) {};}
\node[mpsnode] at (\xD,\ymps) {};
\node[mpofull] at (\xD,\ympo) {};

% ---- the change of basis, drawn as the rank-2 tensor it is ---------------
%  U is a d x D isometry, so it is a rank-2 box sitting ON the physical leg:
%  U^dagger takes the MPS's d-dimensional index up to the full MPO's D, and U
%  brings the output index back down.  Labels sit to the LEFT of the leg; the
%  strip to its right is kept clear for the zoom callout into the inset.
\node[ubox] at (\xD,\yUa) {$U^\dagger$};
\node[ubox] at (\xD,\yUb) {$U$};
\node[figannot,anchor=east] at (\xdim,348.5) {$d$};
\node[figannot,anchor=east] at (\xdim,324)   {$D$};
\node[figannot,anchor=east] at (\xdim,297)   {$D$};
\node[figannot,anchor=east] at (\xdim,272.5) {$d$};
\node[figannot,anchor=north west,align=left,inner sep=0pt] at (\xcol,378)
  {at every site\\[1pt]
   $\textcolor{figblueA}{\widehat{\mathbb{W}}} \to
    U\textcolor{figblueA}{\widehat{\mathbb{W}}}U^\dagger$\\[1pt]
   $D{=}65 \to d{=}32$};

% ===========================================================================
%  2.  the equals sign and the eigenvector, on ONE line.  The two sides ARE
%      the eigenvalue equation: above is W|P_i>, here is mu_i |P_i>.  The
%      symbolic line that used to sit between them has been deleted -- the
%      diagram states it once, and writing it again in symbols is the same
%      claim in two registers.
%
%      The result MPS carries no v_0 ... v_L labels.  Its legs run over the d
%      retained coefficients, not over node voltages; which node each site IS
%      is said by the chain drawn directly beneath it.
% ===========================================================================
\def\yres{235}
% anchored EAST at 100, not west at 78: the first MPS circle reaches x=102
% and is drawn opaquely afterwards, so a west-anchored label loses its
% subscript underneath it.
\node[anchor=base east] at (100,\yres) {$=\ \textcolor{figblueA}{\mu_i}$};
\brokenline{bond}{\rA}{\rB}{\yres}
\draw[bond] (\rB,\yres) -- (\rC,\yres);
\brokenline{bond}{\rC}{\rD}{\yres}
\foreach \x in {\rA,\rB,\rC,\rD}{%
  \draw[bond] (\x,\yres) -- (\x,215);
  \nodedot{\x}{215}}
\foreach \x in {\rA,\rB,\rC,\rD}{\node[mpsnode] at (\x,\yres) {};}
\node[figannot,anchor=south] at (173,238) {$\chi$};
\node[figannot,anchor=west]  at (193,223) {$d$};

% ===========================================================================
%  3.  the CMOS chain, directly beneath the result MPS and on its abscissae.
%      Three cross-coupled inverter pairs, four nodes.  The triangles reach
%      y = \ywire + 20.8 = 205.8, against the MPS free-leg dots at 215.
% ===========================================================================
\def\ywire{185}
%  Stubs of bare wire: 5pt at the two outer nodes, 6pt at the two inner ones,
%  leaving 19pt for each elision.  Unit 2 is centred in the unelided span.
\draw[figwire] (\rA,\ywire) -- (115,\ywire);
\cmosunit{115}{\ywire}
\brokenline{figwire}{137}{\rB}{\ywire}
\draw[figwire] (\rB,\ywire) -- (162,\ywire);
\cmosunit{162}{\ywire}
\draw[figwire] (184,\ywire) -- (\rC,\ywire);
\brokenline{figwire}{\rC}{209}{\ywire}
\cmosunit{209}{\ywire}
\draw[figwire] (231,\ywire) -- (\rD,\ywire);
\foreach \x in {\rA,\rB,\rC,\rD}{\nodedot{\x}{\ywire}}
% the ONLY place v_0 ... v_L appear: on the chain, where they are voltages.
\node[figannot,anchor=north] at (\rA,160) {$v_0$};
\node[figannot,anchor=north] at (\rB,160) {$v_i$};
\node[figannot,anchor=north] at (\rC,160) {$v_{i+1}$};
\node[figannot,anchor=north] at (\rD,160) {$v_L$};

% ===========================================================================
%  4.  zoom callout: the rows of U, i.e. the retained basis it projects onto.
%      Set beside the sandwich, level with it, so the wedge is short and the
%      DMRG arrows below are free to descend from the equation.
% ===========================================================================
\draw[callout] (141,292.5) -- (157,338);
\draw[callout] (141,281.5) -- (157,274);

\begin{axis}[figaxis,
  at={(158pt,272pt)}, anchor=south west,
  width=76pt, height=68pt,
  xlabel={$v_L/v_e$},
  xmin=-33, xmax=33, xtick={-30,0,30},
  ymin=-0.072, ymax=0.072, ytick=\empty,
  xlabel shift=-2pt,
]
% the vertical scale of an orthonormalised basis vector is convention, so the
% y ticks and label carry no information; dropped, and the width recovered.
\draw[figgray,line width=0.4pt] (axis cs:-33,0) -- (axis cs:33,0);
\addplot[figgreen,line width=0.7pt,smooth] table[x=x,y=f4] {fig2_basis.dat};
\addplot[figgreen,line width=0.7pt,smooth] table[x=x,y=f3] {fig2_basis.dat};
\addplot[figgreen,line width=0.7pt,smooth] table[x=x,y=f2] {fig2_basis.dat};
\addplot[figblueB,figemph,smooth]          table[x=x,y=f1] {fig2_basis.dat};
\addplot[figblueA,figemph,smooth]          table[x=x,y=f0] {fig2_basis.dat};
\end{axis}

% ===========================================================================
%  5.  the two DMRG variants, in the left margin the chain leaves free, each
%      landing on the eigenvector it returns and its eigenvalue.  The arrows
%      fan DOWN-LEFT out of the equation: down, so the reader still finishes
%      at the bottom of the figure; left, into the column beside the chain.
%      No sketches of the eigenvectors: for L>1 they are functions of all L+1
%      node voltages, and a cartoon of them invites the reader to think the
%      flip mechanism can be read off the picture.  The inset above already
%      shows real single-unit eigenvectors, the only honest curve here.
% ===========================================================================
%  Both arrows leave from just under the leftmost result-MPS site, clear of
%  its circle (which reaches x=102, y=227) and of the descender of mu_i, and
%  land just outside the upper right corner of the group they name.  Measured
%  label extents at 8pt: "DMRG" spans x=28..67, "Excited State" x=19..77, so
%  the tips have to stop at x>67 and x>79 respectively.
\draw[fatarrow] (101,227) -- (73,211);
\draw[fatarrow] (103,224) -- (84,164);

%  Spacing is 4pt within a group and 10pt between them, so the two eigenpairs
%  read as two groups rather than as six stacked lines.  Everything here is
%  centred on x=48; the widest line, mu_1, spans x=17..79 and so clears the
%  chain, whose leftmost ink is the node dot at x=108.
\node[figannot,anchor=north,figblueA]              at (48,214) {DMRG};
\node[anchor=north]                                at (48,201) {$\ket{P_{\rm ss}}$};
\node[figannot,anchor=north]                       at (48,184) {$\mu_0 = 0$};

\node[figannot,anchor=north,align=center,figblueA] at (48,164)
  {Excited State\\DMRG};
\node[anchor=north]                                at (48,140) {$\ket{P_{\rm relax}}$};
\node[figannot,anchor=north]                       at (48,123)
  {$\mu_1 = 2\langle\tau_{\rm err}\rangle^{-1}$};

\end{tikzpicture}
\end{document}
